\documentclass[12pt]{article}
\usepackage{amsmath,amssymb,amsfonts,amsthm}
\usepackage[margin=1in]{geometry}

\begin{document}

\begin{center}
{\Large\bfseries Chapter 4, Problem 9}\\[6pt]
Byron \& Fuller, \textit{Mathematics of Classical and Quantum Physics}
\end{center}

\bigskip

\textbf{Problem.} Define $\exp(A)$, where $A$ is any matrix, by the power series expansion of $\exp(x) = 1 + x + \cdots$. Show that if there exists a matrix $P$ which reduces $A$ to diagonal form by a similarity transform, then one has
\[
\det\exp(A) = \exp\operatorname{tr}(A).
\]

\bigskip
\textbf{Solution.}

Since $A$ is diagonalizable, there exists an invertible matrix $P$ such that:
\[
A = P D P^{-1}, \quad D = \operatorname{diag}(\lambda_1, \lambda_2, \ldots, \lambda_n),
\]
where $\lambda_i$ are the eigenvalues of $A$.

First, note that $A^k = P D^k P^{-1}$, so:
\[
e^A = \sum_{k=0}^{\infty} \frac{A^k}{k!} = P\left(\sum_{k=0}^{\infty} \frac{D^k}{k!}\right)P^{-1} = P\, e^D\, P^{-1}.
\]

Since $D$ is diagonal, $e^D = \operatorname{diag}(e^{\lambda_1}, e^{\lambda_2}, \ldots, e^{\lambda_n})$.

Now take the determinant:
\[
\det(e^A) = \det(P\, e^D\, P^{-1}) = \det(P)\det(e^D)\det(P^{-1}) = \det(e^D).
\]
Since $e^D$ is diagonal:
\[
\det(e^D) = \prod_{i=1}^n e^{\lambda_i} = e^{\sum_{i=1}^n \lambda_i} = e^{\operatorname{tr}(D)}.
\]

Finally, since the trace is invariant under similarity transformations:
\[
\operatorname{tr}(A) = \operatorname{tr}(P D P^{-1}) = \operatorname{tr}(D) = \sum_{i=1}^n \lambda_i.
\]

Therefore:
\[
\det\exp(A) = e^{\operatorname{tr}(D)} = e^{\operatorname{tr}(A)} = \exp\operatorname{tr}(A). \qed
\]

\end{document}
